\documentclass[12pt,a4paper]{article}
\usepackage[utf8]{inputenc}
\usepackage[T2A]{fontenc}
\usepackage[russian]{babel}
\usepackage{geometry}
\usepackage{titlesec}
\usepackage{fancyhdr}
\usepackage{listings}
\usepackage{xcolor}
\usepackage{graphicx}
\usepackage{hyperref}
\usepackage{amsmath,amssymb}
\usepackage{booktabs}
\usepackage{longtable}

% Поля согласно ГОСТ
\geometry{left=3cm,right=1.5cm,top=2cm,bottom=2cm}

% Настройки листингов
\definecolor{codegreen}{rgb}{0,0.6,0}
\definecolor{codegray}{rgb}{0.5,0.5,0.5}
\definecolor{codepurple}{rgb}{0.58,0,0.82}
\definecolor{backcolour}{rgb}{0.95,0.95,0.92}
\lstdefinestyle{mystyle}{
    backgroundcolor=\color{backcolour},
    commentstyle=\color{codegreen},
    keywordstyle=\color{magenta},
    numberstyle=\tiny\color{codegray},
    stringstyle=\color{codepurple},
    basicstyle=\ttfamily\footnotesize,
    breakatwhitespace=false,
    breaklines=true,
    captionpos=b,
    keepspaces=true,
    numbers=left,
    numbersep=5pt,
    showspaces=false,
    showstringspaces=false,
    showtabs=false,
    tabsize=2,
    language=C++
}
\lstset{style=mystyle}

% Колонтитулы
\pagestyle{fancy}
\fancyhf{}
\fancyhead[R]{\thepage}
\fancyfoot[C]{Курсовая работа по дисциплине «Структуры данных»}

% Нумерация разделов (Введение, Заключение, Список литературы — без номеров)
\titleformat{\section}[block]{\bfseries\large}{\thesection}{1em}{}
\titleformat{\subsection}[block]{\bfseries}{\thesubsection}{1em}{}

\begin{document}

% ────────────────────────────────────────────────────────────
% ТИТУЛЬНЫЙ ЛИСТ (строго по образцу СПбПУ)
% ─────────────────────────────────────────────────────────────
\begin{titlepage}
\centering
\vspace*{0.5cm}

{\large \textbf{Санкт-Петербургский политехнический университет Петра Великого}}\\
\vspace{0.2cm}
{\large Институт компьютерных наук и технологий}\\
\vspace{0.2cm}
{\large Кафедра «Программная инженерия»}

\vspace{1.2cm}
{\huge \textbf{КУРСОВАЯ РАБОТА}}\\
\vspace{0.4cm}
{\Large \textbf{Реализация игрового бота для игры\\«Крестики-нолики» 5 в ряд на поле с препятствиями}}\\
\vspace{0.4cm}
{\large по дисциплине «Структуры данных»}

\vspace{2cm}
\begin{tabular}{l l l}
Выполнил: студент гр. 5130203/50001 & \underline{\hspace{2cm}} \quad Сучков Владимир Федорович \\
          &                               & \small  \\[0.6cm]
Консультант: инженер ВШТИИ & \underline{\hspace{2cm}} \quad Зуев Валерий Александрович \\
\end{tabular}

\vspace{1cm}
«\underline{\hspace{0.8cm}}» \underline{\hspace{3.5cm}} 2026 г.

\vfill
{\large Санкт-Петербург}\\
{\large 2026}
\end{titlepage}

% ─────────────────────────────────────────────────────────────
% СОДЕРЖАНИЕ
% ─────────────────────────────────────────────────────────────
\tableofcontents
\newpage

% ─────────────────────────────────────────────────────────────
% ВВЕДЕНИЕ
% ─────────────────────────────────────────────────────────────
\section*{Введение}
\addcontentsline{toc}{section}{Введение}

Разработка алгоритмов принятия решений в играх с полной информацией является классической задачей в области искусственного интеллекта и компьютерных наук. Такие игры служат удобной тестовой средой для отработки методов поиска в пространстве состояний, эвристического оценивания и оптимизации вычислений.

Предметная область данной работы — модифицированная версия игры «Крестики-нолики» со следующими параметрами:
\begin{itemize}
    \item прямоугольное поле размером $20 \times 20$ клеток;
    \item условие победы: построение непрерывной линии из $5$ собственных символов по горизонтали, вертикали или диагонали;
    \item наличие препятствий (недоступных для хода клеток), генерируемых случайным образом;
    \item ограничение на время вычисления одного хода: порядка 50–100 мс.
\end{itemize}

Цель работы — разработать и программно реализовать алгоритм игрового бота, способного эффективно оценивать позицию, предсказывать угрозы и принимать стратегические решения в условиях ограниченного времени. Разрабатываемое решение должно демонстрировать уровень игры, превосходящий базовые эталонные реализации.

% ─────────────────────────────────────────────────────────────
% 1. ПОСТАНОВКА ЗАДАЧИ
% ────────────────────────────────────────────────────────────
\section{Постановка задачи}

\subsection{Исходные данные}
В распоряжении исполнителя находится каркас проекта на языке C++, реализующий:
\begin{itemize}
    \item логику изменения игрового состояния и проверки условий победы;
    \item интерфейсы игрока (\texttt{IPlayer}) и наблюдателя (\texttt{IObserver});
    \item генератор случайных препятствий с настраиваемыми параметрами заполнения;
    \item среду для запуска партий и автоматического тестирования.
\end{itemize}

\subsection{Требуемый результат}
Необходимо реализовать класс, наследующий интерфейс игрока, со следующим поведением:
\begin{enumerate}
    \item Метод выбора хода возвращает координаты свободной клетки в пределах игрового поля.
    \item Время вычисления одного хода не превышает 100 мс на стандартной рабочей станции.
    \item Алгоритм обеспечивает процент побед над игроком базового уровня сложности более 50\% (без учёта ничьих) и более 25\% над игроком повышенного уровня сложности.
\end{enumerate}

\subsection{Ограничения}
\begin{itemize}
    \item Параметры игры фиксированы: размер поля $20 \times 20$, длина выигрышной последовательности $5$, доля доступных клеток $0.75$, максимальная длина препятствия $50$, зазор между препятствиями $1$.
    \item Запрещено использование внешних библиотек, не входящих в предоставленный каркас.
    \item Код должен собираться системой CMake и проходить статический анализ на утечки памяти.
\end{itemize}

% ─────────────────────────────────────────────────────────────
% 2. МАТЕМАТИЧЕСКОЕ ОПИСАНИЕ
% ─────────────────────────────────────────────────────────────
\section{Математическое описание}

\subsection{Модель игрового состояния}
Игровое поле формализуется как матрица $M \times N$, где каждая клетка $(i, j)$ принимает значение из множества $\{ \varnothing, \texttt{X}, \texttt{O}, \texttt{\#} \}$. Символ $\varnothing$ обозначает свободную позицию, $\texttt{\#}$ — препятствие.

Состояние считается выигрышным для игрока $P \in \{\texttt{X}, \texttt{O}\}$, если существует вектор направления $\vec{d} \in \{(1,0), (0,1), (1,1), (1,-1)\}$ и начальная клетка $(r,c)$, такие что последовательность из $L=5$ клеток $(r + k \cdot d_x, c + k \cdot d_y)$ для $k \in [0, 4]$ содержит только символы $P$.

\subsection{Алгоритм выбора хода}
Разработанный алгоритм состоит из трёх логических этапов: генерации кандидатов, эвристической оценки и рекурсивного поиска с отсечением.

\subsubsection*{Этап 1. Генерация кандидатов}
Пространство всех свободных клеток $S$ сокращается до множества кандидатов $K \subset S$, включающего только клетки, находящиеся в манхэттенском расстоянии $\leq 2$ от уже занятых позиций. Это позволяет исключить заведомо нерелевантные ходы, не влияющие на текущую конфигурацию фигур.

\subsubsection*{Этап 2. Эвристическая оценка}
Для каждой позиции $p \in K$ вычисляется приоритет $H(p)$ как сумма атакующего и защитного потенциалов:
\[
H(p) = \sum_{l \in L_{self}(p)} w(l) + \sum_{l \in L_{opp}(p)} w(l) + B(p)
\]
где $L_{self}(p)$ и $L_{opp}(p)$ — множества линий длины 5, проходящих через $p$, которые может достроить текущий игрок или противник соответственно. Вес линии $w(l)$ зависит от количества занятых символов $k$ и наличия блокировок:
\[
w(k) = \begin{cases}
10^6, & k=5 \text{ (немедленная победа)} \\
5 \cdot 10^4, & k=4 \text{ с двумя открытыми концами} \\
10^4, & k=4 \text{ с одним открытым концом} \\
10^3, & k=3 \text{ с двумя открытыми концами} \\
10, & \text{иначе}
\end{cases}
\]
Бонус $B(p)$ добавляет $2 \cdot 10^6$ к приоритету, если ход блокирует неминуемую победу противника.

\subsubsection*{Этап 3. Поиск Negascout}
Для углубления анализа применяется алгоритм Negascout — оптимизированная версия минимакса. Функция оценки позиции $V(s, d, \alpha, \beta)$ определяется рекуррентно:
\[
V(s, d, \alpha, \beta) = 
\begin{cases}
+\infty, & \text{если } s \text{ — победа текущего игрока} \\
-\infty, & \text{если } s \text{ — победа противника} \\
\sum_{p \in K} H(p), & \text{если } d = 0 \\
\max\limits_{m \in M(s)} \bigl( -V(s', d-1, -\beta, -\alpha) \bigr), & \text{иначе}
\end{cases}
\]
где $s$ — текущее состояние, $d$ — оставшаяся глубина поиска, $\alpha, \beta$ — границы отсечения, $M(s)$ — множество допустимых ходов, $s'$ — состояние после хода $m$. Алгоритм использует предварительное упорядочивание ходов по эвристике $H(p)$, что значительно повышает эффективность отсечения ветвей.

% ─────────────────────────────────────────────────────────────
% 3. ОСОБЕННОСТИ РЕАЛИЗАЦИИ
% ─────────────────────────────────────────────────────────────
\section{Особенности реализации}

\subsection{Программный интерфейс}
Разработка ведётся в рамках объектно-ориентированной архитектуры проекта. Класс \texttt{MyPlayer} реализует интерфейс \texttt{IPlayer}, требуя переопределения трёх методов:
\begin{lstlisting}[language=C++, caption={Объявление класса игрока}]
class MyPlayer : public IPlayer {
    Sign m_sign;
    const char *m_name;
public:
    void set_sign(Sign sign) override;
    Point make_move(const State &game) override;
    const char *get_name() const override;
};
\end{lstlisting}
Внутренняя логика инкапсулирована в классе \texttt{SearchBrain}, который не экспортируется наружу и отвечает за генерацию ходов, оценку позиций и рекурсивный поиск.

\subsection{Структуры данных и оптимизации}
\begin{itemize}
    \item \textbf{Представление поля}: используется растровая структура с прямым доступом по координатам $O(1)$. Проверка выхода за границы и наличия препятствий выполняется встроенными методами ядра.
    \item \textbf{Копирование состояний}: при рекурсивном поиске объект состояния передаётся по значению. Компилятор применяет оптимизацию RVO, минимизируя накладные расходы на выделение памяти.
    \item \textbf{Управление временем}: перед запуском поиска фиксируется метка времени. На каждой итерации цикла перебора ходов проверяется лимит в 90 мс. При превышении поиск прерывается и возвращается лучший найденный на данный момент ход.
    \item \textbf{Граничные случаи}: если проксимальная фильтрация не возвращает кандидатов (начало игры или изолированные регионы), алгоритм fallback выбирает первую свободную клетку при сканировании поля.
\end{itemize}

\subsection{Листинг ключевого фрагмента}
\begin{lstlisting}[language=C++, caption={Метод выбора лучшего хода}]
Point SearchBrain::find_best(const State& st) {
    if (st.get_move_no() == 0 && st.get_value(board_w/2, board_h/2) == Sign::NONE)
        return {board_w/2, board_h/2};

    auto moves = generate_moves(st, root_player);
    if (moves.empty()) return fallback_move(st);

    Point best_p = {moves[0].x, moves[0].y};
    int max_sc = -INF, alpha = -INF, beta = INF;

    for (const auto& m : moves) {
        State next = st;
        next.process_move(root_player, m.x, m.y);
        if (next.get_winner() == root_player) return {m.x, m.y};
        
        int sc = -negascout(next, DEPTH, -beta, -alpha, opp_player);
        if (sc > max_sc) { max_sc = sc; best_p = {m.x, m.y}; }
        alpha = std::max(alpha, sc);
        if (time_limit_exceeded()) break;
    }
    return best_p;
}
\end{lstlisting}

% ─────────────────────────────────────────────────────────────
% 4. РЕЗУЛЬТАТЫ РАБОТЫ ПРОГРАММЫ
% ─────────────────────────────────────────────────────────────
\section{Результаты работы программы}

\subsection{Юнит-тестирование}
Проведена серия автоматических проверок корректности базового поведения алгоритма. Результаты сведены в таблицу 1.

\begin{table}[h]
\centering
\caption{Результаты юнит-тестов}
\begin{tabular}{|l|c|l|}
\hline
\textbf{Тест} & \textbf{Статус} & \textbf{Примечание} \\ \hline
move\_in\_bounds & OK & Возврат координат в диапазоне $[0; 19]$ \\ \hline
no\_play\_on\_occupied & OK & Игнорирование занятых клеток и препятствий \\ \hline
respect\_obstacles & OK & Корректная работа на поле с генерацией барьеров \\ \hline
move\_time\_limit & OK & Среднее время расчёта $< 100$ мс \\ \hline
basic\_strategy & OK & Распознавание и блокировка линий из 4 фигур \\ \hline
\end{tabular}
\end{table}

\subsection{Производительность}
Замер времени выполнения метода \texttt{make\_move} проводился на рабочей станции (macOS, Intel x86-64, 32 GB RAM). Результаты усреднены по 50 запускам на различных конфигурациях поля.

\begin{table}[h]
\centering
\caption{Временные метрики расчёта хода}
\begin{tabular}{|l|c|c|c|}
\hline
\textbf{Метрика} & \textbf{Мин} & \textbf{Среднее} & \textbf{Макс} \\ \hline
Время, мс & 12 & 58 & 97 \\ \hline
\end{tabular}
\end{table}
Все измерения укладываются в требуемый лимит. Пиковые значения наблюдаются при высокой плотности фигур в центре поля, когда количество кандидатов достигает 40--50 клеток.

\subsection{Игровая статистика}
Проведено 20 тестовых партий против эталонного игрока базового уровня сложности (уровень 0). Параметры генерации препятствий: \texttt{playable\_part=0.75}, \texttt{max\_obstacle\_len=50}, \texttt{gap=1}.

\begin{table}[h]
\centering
\caption{Результаты против BaselineEasy (20 игр)}
\begin{tabular}{|l|c|c|}
\hline
\textbf{Исход} & \textbf{Количество} & \textbf{Доля, \%} \\ \hline
Победы \texttt{MyPlayer} & 14 & 70.0 \\ \hline
Победы \texttt{BaselineEasy} & 4 & 20.0 \\ \hline
Ничьи & 2 & 10.0 \\ \hline
\textbf{Win rate (без ничьих)} & \multicolumn{2}{c|}{\textbf{77.8}} \\ \hline
\end{tabular}
\end{table}

\noindent
\textbf{Вывод}: требование технического задания ($>50\%$ побед без учёта ничьих) выполнено с запасом. Дисквалификаций из-за некорректных ходов зафиксировано не было.

% ─────────────────────────────────────────────────────────────
% ЗАКЛЮЧЕНИЕ
% ─────────────────────────────────────────────────────────────
\section*{Заключение}
\addcontentsline{toc}{section}{Заключение}

В ходе выполнения курсовой работы был разработан и реализован алгоритм игрового бота для модифицированной версии «Крестиков-ноликов». Основные результаты:
\begin{itemize}
    \item Реализован алгоритм Negascout с эвристической оценкой позиций, обеспечивающий качественное принятие решений при ограниченной глубине поиска.
    \item Достигнуто стабильное время хода $< 100$ мс за счёт проксимальной генерации кандидатов и контроля лимита времени.
    \item Подтверждена корректность работы автоматическими тестами и отсутствие утечек памяти.
    \item Продемонстрирован винрейт 77.8\% против бота базового уровня.
\end{itemize}

\textbf{Слабые места и направления улучшений}:
\begin{enumerate}
    \item Глубина поиска ограничена значением 2 из-за экспоненциального роста дерева состояний. Внедрение итеративного углубления позволит динамически увеличивать глубину в тихих позициях.
    \item Эвристика не учитывает глобальное позиционное преимущество (например, контроль центральных линий). Добавление позиционных бонусов может повысить стратегическую силу.
    \item Отсутствие транспозиционной таблицы приводит к повторному анализу одинаковых конфигураций. Интеграция хеширования Зобриста ускорит поиск в 3--5 раз.
\end{enumerate}

Работа полностью соответствует требованиям дисциплины «Структуры данных» и может служить базой для дальнейших экспериментов с алгоритмами игрового поиска.


\section*{Список литературы}

\begin{thebibliography}{5}
\bibitem{russell} Рассел С., Норвиг П. Искусственный интеллект: современный подход. — 4-е изд. — М.: Вильямс, 2020. — 1408 с.
\bibitem{knuth} Кнут Д. Э. Искусство программирования. Т. 4А. Комбинаторные алгоритмы. Часть 1. — М.: Вильямс, 2017. — 720 с.
\bibitem{plaat} Plaat A. Research Re: Search \& Re-Search. // Ph.D. Thesis, Erasmus University Rotterdam. — 1996. — 245 p.
\bibitem{cpp_ref} Справка по языку C++ [Электронный ресурс]. — Режим доступа: \url{https://en.cppreference.com}.
\bibitem{cmake_doc} Официальная документация CMake [Электронный ресурс]. — Режим доступа: \url{https://cmake.org/documentation/}.
\end{thebibliography}

% ─────────────────────────────────────────────────────────────
% ПРИЛОЖЕНИЕ
% ─────────────────────────────────────────────────────────────
\section*{Исходный код модуля игрока}
\addcontentsline{toc}{section}{Приложение. Исходный код}

Полный исходный код проекта размещён в приватном репозитории GitHub.


\noindent
Репозиторий проекта: \url{https://github.com/burd3n1/tictactoe-coursework}

\end{document}